\documentclass[11pt]{article}
\usepackage[margin=2.4cm]{geometry}
\usepackage{amsmath, amssymb, amsthm, mathtools}
\usepackage{xcolor}
\usepackage{microtype}
\usepackage{tikz-cd}
\usepackage{booktabs}
\usepackage{enumitem}
\usepackage[hidelinks]{hyperref}

% Math operator spacing for the hat-decorated combinators
\newcommand{\hexp}{\mathop{\widehat{\exp}}}
\newcommand{\hlog}{\mathop{\widehat{\log}}}
\newcommand{\hsqrt}{\mathop{\widehat{\sqrt{\vphantom{x}}}}}
\newcommand{\hsq}{\mathop{\widehat{\mathrm{sq}}}}
\newcommand{\hinv}{\mathop{\widehat{\mathrm{inv}}}}
\newcommand{\hadd}{\mathbin{\widehat{+}}}
\newcommand{\hsub}{\mathbin{\widehat{-}}}
\newcommand{\hmul}{\mathbin{\widehat{\cdot}}}
\newcommand{\hdiv}{\mathbin{\widehat{/}}}
\newcommand{\hpow}{\mathbin{\widehat{\hat{}}\,}}
% Complex variants
\newcommand{\hexpC}{\mathop{\widehat{\exp_{\!\mathbb{C}}}}}
\newcommand{\hlogC}{\mathop{\widehat{\log_{\!\mathbb{C}}}}}
\newcommand{\haddC}{\mathbin{\widehat{+_{\!\mathbb{C}}}}}
\newcommand{\hsubC}{\mathbin{\widehat{-_{\!\mathbb{C}}}}}
\newcommand{\hmulC}{\mathbin{\widehat{\cdot_{\!\mathbb{C}}}}}
\newcommand{\hdivC}{\mathbin{\widehat{/_{\!\mathbb{C}}}}}

\newcommand{\F}{\mathsf{F36}}
\newcommand{\EL}{\mathsf{EL}}
\newcommand{\EML}{\mathsf{EML}}
\newcommand{\EMLC}{\mathsf{EML}_{\mathbb{C}}}
\newcommand{\Eset}{\mathcal{E}}
\newcommand{\EsetC}{\mathcal{E}_{\mathbb{C}}}

\theoremstyle{plain}
\newtheorem{theorem}{Theorem}
\newtheorem{lemma}[theorem]{Lemma}
\newtheorem{proposition}[theorem]{Proposition}

\theoremstyle{definition}
\newtheorem{definition}[theorem]{Definition}
\newtheorem{remark}[theorem]{Remark}
\newtheorem{example}[theorem]{Example}

\title{Structure of the \textsc{eml} completeness proof\\[2pt]
       \large An expository tour of grammars, combinators, and bridges}
\author{}
\date{}

\begin{document}
\maketitle

\begin{abstract}
We give a high-level account of how the \emph{completeness theorem}
\textemdash\ \emph{every elementary function is expressible as a finite
expression in the binary operator $\mathrm{eml}(a, b) = \exp(a) -
\log(b)$, paired with the constant $1$} \textemdash\ is proved.
The proof rests on a layered architecture: three inductive grammars
($\F$, $\EL$, $\EML$) connected by two compilers; a catalogue of
\emph{derived combinators} that recover the usual repertoire
($+$, $-$, $\cdot$, $/$, $\exp$, $\log$, $\sqrt{\,}$, $\ldots$) as
fixed-shape sub-trees of the $\mathrm{eml}$ syntax; a single structural
correctness theorem for the compiler; and a small zoo of
\emph{Euler bridges} from the complex EML grammar back to the real
elementary functions. We trace one paper claim end-to-end as a worked
example, and finish with the scoreboard summarising what is sealed.
A closing addendum (\S\ref{sec:post-submission}) records the post-
submission widening to full real-domain trig (Path~C$'$), the cleanup
of the Sheffer companion scaffolding (Plan~A), and partial progress
on per-primitive completeness for the paper's Sheffer cousins
(Plans~D and~E).
\end{abstract}

\tableofcontents

% =====================================================================
\section{The headline theorem}
% =====================================================================

The paper's claim, stripped to bare mathematical content, is the
following.

\begin{theorem}[Headline]
\label{thm:headline}
Let $\F$ denote the catalogue of $36$ paper primitives \textemdash\ atoms
($\mathbf{1},\,\pi,\,e,\,-1,\,2,\,\tfrac{1}{2}$ and the input variable),
the eight standard real unaries, the six hyperbolic functions, the
six trigonometric and inverse-trigonometric functions, and the eight
binary operations. Let $\Eset$ be the single-binary-operator grammar
\[
  \Eset \;\ni\; T \;::=\; \mathbf{1} \;\;\big|\;\; \mathbf{x}_n
       \;\;\big|\;\; \mathrm{eml}(T, T),
  \qquad \mathrm{eml}(a, b) := \exp(a) - \log(b).
\]
Then for each $f \in \F$ there is a finite expression $\mathsf{T}_{f}
\in \Eset$ \textup{(}or in its complex variant $\EsetC$\textup{)} such
that on a non-empty open subset of $f$'s natural domain, the
partial evaluation of $\mathsf{T}_f$ yields a value whose appropriate
projection $\Re$ or $\Im$ equals $f(x)$.
\end{theorem}

The hedge \emph{non-empty open subset} is honest. Three measure-zero
points \textemdash\ $\sqrt{0}$, $\mathrm{arcosh}(1)$, $\mathrm{hypot}(0,0)$
\textemdash\ lie outside the natural construction; a few trig primitives
are sealed on strict-but-symmetric subdomains rather than on the full
real line. The bulk of the result, however, is a clean structural
completeness, and the architecture below explains how the bulk is
obtained.

The aim of this report is to explain \emph{why} Theorem
\ref{thm:headline} is true: what shape the proof has, what each
architectural layer contributes, and where the interesting ideas live.

% =====================================================================
\section{Three languages, two compilers}
% =====================================================================

The proof reduces the problem $f \rightsquigarrow \mathsf{T}_f$ through
a three-language pipeline. The picture is

\begin{center}
\begin{tikzcd}[column sep=large, row sep=small]
\boxed{\F}
  \arrow[r, "\text{translate?}"]
  \arrow[rrd, dashed, bend right=12, "\text{paper claim}"' inner sep=1pt]
& \boxed{\EL}
  \arrow[r, "\text{compile}"]
& \boxed{\EML}
  \arrow[d, "\iota"]
\\[-0.5em]
& & \boxed{\EMLC}
\end{tikzcd}
\end{center}

The four boxes are inductive types; the three arrows are total
functions; the dashed arrow names the public statement we want to
verify.

\paragraph{$\F$.} The paper's source language. It has $36$ named
constructors and a partial denotation: each primitive returns $\bot$
outside its natural mathematical domain (e.g.\ $\log$ on $\le 0$).

\paragraph{$\EL$.} An intermediate language whose alphabet is
\emph{elementary}: atoms ($\mathbf{1},\,\mathbf{x}_n$, named constants),
the unaries $\{\exp,\log,\mathrm{neg},\mathrm{halve},\mathrm{sq},\sqrt{\,},
\mathrm{inv}\}$, and the binaries $\{+,-,\cdot,/,\hat{},\log_b,
\mathrm{avg},\mathrm{hypot}\}$. The translation $\F \to \EL$ expands
non-elementary names through standard identities: e.g.
\begin{align*}
  \sinh(a) &\mapsto \tfrac{1}{2}\bigl(\exp(a) - \exp(-a)\bigr), &
  \tanh(a) &\mapsto \sinh(a) / \cosh(a), \\
  \mathrm{arsinh}(a) &\mapsto \log\bigl(a + \sqrt{a^2 + 1}\bigr), &
  \mathrm{artanh}(a) &\mapsto \tfrac{1}{2}\bigl(\log(1 + a) - \log(1 - a)\bigr).
\end{align*}
The language $\EL$ is small enough to compile, but rich enough to
absorb every named primitive of $\F$.

\paragraph{$\EML$.} The pure single-operator grammar of
Theorem~\ref{thm:headline}. The compiler $\EL \to \EML$ is the heart of
the proof.

\paragraph{$\EMLC$.} The complex-coefficient version of the EML
grammar, with the same syntax but complex evaluation. Trig and the
constants $\pi,\,i$ require Euler's formula and therefore complex
intermediates; $\EMLC$ is where they live. The embedding $\iota$ is a
\emph{homomorphism}: it does nothing on syntax (every constructor
maps to itself) and changes only the evaluation rule.

\medskip

The architectural payoff is two short statements.

\begin{theorem}[Real-fragment compiler]
\label{thm:real-compiler}
Suppose $f \in \F$ satisfies $\mathrm{translate?}(f) = \mathrm{some}\,e$
with $e \in \EL$. Then for every environment
$\mathrm{env}:\mathbb{N}\to\mathbb{R}$ and every value $v\in\mathbb{R}$,
\[
  f.\mathrm{eval?}(\mathrm{env}) = \mathrm{some}\,v
  \;\Longrightarrow\;
  \mathrm{compile}(e).\mathrm{eval?}(\mathrm{env}) = \mathrm{some}\,v.
\]
\end{theorem}

\begin{theorem}[Complex bridges]
\label{thm:complex-bridges}
For each primitive in $\{\pi,\,i,\,\cos,\,\sin,\,\tan,\,\arctan,\,\arccos,\,\arcsin\}$,
there is a finite witness $\mathsf{C}_f \in \EsetC$ and a projection
$\Pi \in \{\Re,\Im\}$ such that, on a stated open subdomain of $f$'s
natural domain,
\[
  \mathsf{C}_f \downarrow v
  \quad\text{with}\quad
  \Pi(v) = f(x).
\]
\end{theorem}

The remaining sections open each box.

% =====================================================================
\section{The kernel: \texorpdfstring{$\EML$}{EML} grammar and partial evaluation}
% =====================================================================

\begin{definition}[$\EML$ grammar]
\label{def:grammar}
The set $\Eset$ of EML expressions is the smallest set generated by
\[
  \Eset \;\ni\; T \;::=\; \mathbf{1} \;\big|\; \mathbf{x}_n \;\big|\;
   \mathrm{eml}(T,T),
\]
where $n\in\mathbb{N}$. The complex variant $\EsetC$ has the same
syntax; only the evaluation rule changes.
\end{definition}

\begin{definition}[Partial evaluation]
\label{def:eval}
Fix an environment $\mathrm{env}: \mathbb{N} \to \mathbb{C}$. Define a
partial function $\downarrow$ on $\EsetC$ by
\[
  \mathbf{1} \downarrow 1, \qquad
  \mathbf{x}_n \downarrow \mathrm{env}(n),
\]
\[
  \mathrm{eml}(a, b) \downarrow \exp(\alpha) - \log(\beta)
  \quad\text{when}\quad a \downarrow \alpha,\;
  b \downarrow \beta,\; \beta \ne 0,
\]
and undefined otherwise. The same rule with $\mathrm{env}:\mathbb{N}\to\mathbb{R}$
defines the real-evaluation variant on $\Eset$, where additionally the
side condition $\beta > 0$ is required (so that $\log\beta\in\mathbb{R}$).
\end{definition}

\paragraph{Why partial.}
A total evaluation would force a junk value at the boundary, e.g.\ the
convention $\log 0 = 0$. Bridge theorems would then have to dance
around the points where this convention disagrees with the natural
mathematical answer. Partial evaluation lets us state, cleanly, "the
witness has no value at the boundary, and the paper claim is stated
on a subset that excludes it" \textemdash\ which is also what the paper
does.

% =====================================================================
\section{The combinator algebra}
\label{sec:combinators}
% =====================================================================

The single binary operator $\mathrm{eml}$ is too low-level to write
useful witnesses with directly. The proof accumulates a catalogue of
\emph{derived combinators}, each defined as a fixed-shape sub-tree of
$\Eset$ (or $\EsetC$), each accompanied by a \emph{closure lemma}
stating the conditions under which the combinator evaluates correctly.

\begin{definition}[Catalogue]
\label{def:catalogue}
The catalogue contains the following derived combinators. Each row
gives the name, the value computed under the closure conditions, and
the size of the underlying $\Eset$-tree.

\smallskip
\begin{center}
\begin{tabular}{l l r}
\toprule
\textbf{Combinator} & \textbf{Computes (with appropriate side conditions)} & \textbf{Tree size} \\
\midrule
$\hexp(T)$       & $\exp(\alpha)$ & $2$ \\
$\hlog(T)$       & $\log(\alpha)$, principal branch & $4$ \\
$A \hadd B$      & $\alpha + \beta$ & $11$ \\
$A \hsub B$      & $\alpha - \beta$ & $7$ \\
$A \hmul B$      & $\alpha \cdot \beta$ via $\exp(\log\alpha + \log\beta)$ & $\approx 18$ \\
$A \hdiv B$      & $\alpha / \beta$ & $\approx 22$ \\
$\hsqrt(T)$      & $\sqrt{\alpha}$ via $\exp(\tfrac{1}{2}\log\alpha)$ & $\approx 10$ \\
$\hsq(T)$        & $\alpha^{2}$ & $\approx 14$ \\
$\hinv(T)$       & $\alpha^{-1}$ via $\exp(-\log\alpha)$ & $\approx 8$ \\
$T \hpow U$      & $\alpha^{\beta}$ via $\exp(\beta \log\alpha)$ & $\approx 28$ \\
\bottomrule
\end{tabular}
\end{center}
\smallskip

The complex variants $\hexpC,\hlogC,\haddC,\hsubC,\hmulC,\hdivC$ have
the same shape; their closure lemmas include extra side conditions on
imaginary parts (the principal-branch strip).
\end{definition}

\begin{remark}[The shape of a closure lemma]
\label{rmk:closure}
Every closure lemma has the form
\[
  \begin{aligned}
   \text{premise:}\quad & A \downarrow \alpha,\; B \downarrow \beta;\\
   \text{side conditions:}\quad & \text{a fixed list of inequalities and}\\
   & \text{non-vanishing conditions on $\alpha,\beta$;}\\
   \text{conclusion:}\quad & A \,\widehat{\star}\, B \downarrow \alpha \star \beta.
  \end{aligned}
\]
The proof is mechanical: unfold the $\mathrm{eml}$-tree from leaves
outward, applying Definition~\ref{def:eval} at each node and threading
the side conditions to ensure that every $\log\circ\exp$ or
$\exp\circ\log$ inverts. A tree of depth $d$ produces a closure
proof with roughly $d$ "step" lines, each invoking one rule.
\end{remark}

\paragraph{The addition combinator, made explicit.}
For concreteness, here is the smallest non-obvious case. Define

\begin{align*}
  L_A &:= \mathrm{eml}\!\bigl(\mathbf{1},\,
   \mathrm{eml}(\mathrm{eml}(\mathbf{1},\mathrm{eml}(A,\mathbf{1})),\mathbf{1})\bigr), \\[2pt]
  R_{A,B} &:= \mathrm{eml}\!\bigl(
   \mathrm{eml}(\mathrm{eml}(\mathbf{1},\mathrm{eml}(A,\mathrm{eml}(A,\mathbf{1}))),\mathbf{1}),\,
   \mathrm{eml}(B,\mathbf{1})\bigr), \\[2pt]
  A \hadd B &:= \mathrm{eml}\!\bigl(L_A,\,\mathrm{eml}(R_{A,B},\mathbf{1})\bigr).
\end{align*}

The closure proof unfolds $A \hadd B$ in eleven sequential steps,
exhibiting the chain
\[
  \begin{aligned}
   \mathrm{eml}(A,\mathbf{1}) &\downarrow e^{\alpha}, &
   \mathrm{eml}(\mathbf{1},\mathrm{eml}(A,\mathbf{1})) &\downarrow e - \alpha, &
   L_A &\downarrow \alpha,\\
   \mathrm{eml}(A,\mathrm{eml}(A,\mathbf{1})) &\downarrow e^{\alpha} - \alpha, &
   \mathrm{eml}(B,\mathbf{1}) &\downarrow e^{\beta}, &
   R_{A,B} &\downarrow e^{\alpha} - \alpha - \beta,\\
  \end{aligned}
\]
\[
  \mathrm{eml}(R_{A,B},\mathbf{1})\downarrow e^{e^{\alpha}-\alpha-\beta},
  \qquad
  A \hadd B \downarrow e^{\alpha} - \bigl(e^{\alpha}-\alpha-\beta\bigr) = \alpha + \beta.
\]
The arithmetic identity at the heart of the construction is just
$\alpha = e - \log(\exp(e - \alpha))$, the principal-branch involution
applied twice.

% =====================================================================
\section{The real-fragment compiler}
\label{sec:real-compiler}
% =====================================================================

With the combinator catalogue in hand, the compiler $\EL \to \EML$ is
straightforward: a structural recursion on $\EL$ expressions in which
each constructor is sent to its catalogue counterpart.

\begin{definition}[Compiler]
\label{def:compile}
The function $\mathrm{compile} : \EL \to \EML$ is defined by recursion.
On atoms,
\[
  \mathrm{compile}(\mathbf{1}) := \mathbf{1}, \qquad
  \mathrm{compile}(\mathbf{x}_n) := \mathbf{x}_n, \qquad
  \mathrm{compile}(\mathbf{0}) := \mathsf{Z};
\]
on unary constructors,
\[
  \mathrm{compile}(\exp e) := \hexp(\mathrm{compile}(e)), \qquad
  \mathrm{compile}(\log e) := \hlog(\mathrm{compile}(e)),
\]
and similarly for the other unaries; on binaries, writing
$\widetilde{e_i} := \mathrm{compile}(e_i)$ for brevity,
\[
  \mathrm{compile}(e_1 + e_2) := \widetilde{e_1} \hadd \widetilde{e_2},
  \qquad
  \mathrm{compile}(e_1 \cdot e_2) := \widetilde{e_1} \hmul \widetilde{e_2},
\]
and similarly for $-$, $/$, $\hat{}\,$, $\log_b$, $\mathrm{avg}$,
$\mathrm{hypot}$. Here $\mathsf{Z}$ is a fixed closed witness for $0$
\textemdash\ a small EML tree that evaluates to $0$ unconditionally.
\end{definition}

\begin{theorem}[Compiler correctness]
\label{thm:compile}
For every $e \in \EL$, every $\mathrm{env}:\mathbb{N}\to\mathbb{R}$,
and every value $v \in \mathbb{R}$,
\[
  e.\mathrm{eval?}(\mathrm{env}) = \mathrm{some}\,v
  \;\Longrightarrow\;
  \mathrm{compile}(e).\mathrm{eval?}(\mathrm{env}) = \mathrm{some}\,v.
\]
\end{theorem}

\begin{proof}[Sketch of the structure]
By induction on $e$. Each constructor case dispatches to one closure
lemma from \S\ref{sec:combinators}: the induction hypothesis supplies
the "$A\downarrow\alpha$" premise, the $\EL$-level evaluation side
conditions (e.g.\ $\alpha > 0$ for $\log$, $\beta \ne 0$ for division)
discharge the combinator's side conditions, and the conclusion is
exactly what the closure lemma delivers.
\end{proof}

The architectural point is that the compiler's correctness proof
contains \emph{one combinator-application per constructor case}. All
the heavy machinery (the eleven-step unfolding for $\hadd$, the deep
unfolding for $\hsqrt$, etc.) is amortised inside the catalogue,
proved once and then invoked.

% =====================================================================
\section{The complex extension and Euler bridges}
\label{sec:euler}
% =====================================================================

\paragraph{Why complex.}
Six paper primitives \textemdash\ $\pi$, $i$, $\cos$, $\sin$, $\tan$,
the inverse trig family \textemdash\ cannot be expressed in real EML alone.
The underlying analytic identities (Euler's $e^{i\pi} = -1$, the
exponential--trigonometric conversions) demand complex intermediates.
We extend the grammar to $\EsetC$ \textemdash\ same syntax, complex
evaluation \textemdash\ and copy the combinator catalogue with branch-cut
side conditions.

\paragraph{Closed constants as building blocks.}
Five \emph{closed} expressions appear repeatedly across the trig
witnesses:
\[
  \mathsf{Z} \downarrow 0,\qquad
  \mathsf{C}_2 \downarrow 2,\qquad
  \mathsf{N}_i \downarrow -i,\qquad
  \mathsf{I} \downarrow i,\qquad
  \mathsf{P} \downarrow \pi.
\]
Each is a finite $\EsetC$ tree paired with an unconditional evaluation
lemma. Trig witnesses use them as plug-in factors rather than
re-deriving Euler's formula at every step.

\paragraph{Real-to-complex lift.}
The embedding $\iota : \EML \to \EMLC$ is a structural homomorphism:
each constructor maps to itself, and only the evaluation rule changes.
Its bridge lemma states: if $T.\mathrm{eval?}(\mathrm{env}) =
\mathrm{some}\,v$ in the real semantics, then
$\iota(T).\mathrm{eval?}(\mathrm{env}\!\otimes\!\mathbb{C}) =
\mathrm{some}\,(v : \mathbb{C})$. This lets every complex trig witness
import the \emph{already-sealed} real arithmetic compositions
($\sqrt{1-x^2}$, $x^2 + 1$, $\tfrac{1}{2}$, etc.) without redoing
branch-cut work.

\paragraph{The Euler bridges.}
Each trig witness encodes one Euler identity, then projects to the
real or imaginary part. The full list:
\begin{equation}
\label{eq:euler}
\begin{aligned}
  \cos x &= \Re\bigl(\exp(ix)\bigr), \\
  \sin x &= \Re\bigl(\exp(i(x - \tfrac{\pi}{2}))\bigr), \\
  \tan x &= \Im\!\left(\frac{\exp(2ix) - 1}{1 + \exp(2ix)}\right) \quad \text{(Cayley quotient)}, \\
  \arctan x &= \Im\bigl(\log(1 + ix)\bigr), \\
  \arccos x &= \Im\bigl(\log(x + i\sqrt{1-x^2})\bigr), \\
  \arcsin x &= \Im\bigl(\log(\sqrt{1-x^2} + ix)\bigr).
\end{aligned}
\end{equation}

The right-hand side of each line is a finite $\EsetC$ witness once the
catalogue combinators and the closed constants $\mathsf{I}, \mathsf{P},
\ldots$ are unfolded. The projection clause is then either $\Re$ or
$\Im$, depending on which one carries the target real value.

% =====================================================================
\section{Coverage and limits}
% =====================================================================

\paragraph{The 36 paper primitives, by family.}
\begin{itemize}[topsep=2pt, itemsep=2pt]
  \item \textbf{Atoms (7).} Direct closed witnesses for $\mathbf{1}$,
    $\mathbf{x}_n$, and the named constants.
  \item \textbf{Real unaries (8) and binaries (8).} Sealed via
    Theorem~\ref{thm:compile}: the $\F\to\EL$ translation expands
    each primitive, the compiler then assembles a witness.
  \item \textbf{Hyperbolic family (6).} Same mechanism: $\sinh$,
    $\cosh$, $\tanh$, $\mathrm{arsinh}$, $\mathrm{arcosh}$,
    $\mathrm{artanh}$ each translate to an $\EL$ composition of
    $\exp$, $\log$, $\sqrt{\,}$, and the four arithmetic operations.
  \item \textbf{Trig family (6).} Sealed via the Euler
    bridges~\eqref{eq:euler}, plus the closed-constant scaffolding of
    \S\ref{sec:euler}.
\end{itemize}

\paragraph{Three structural boundary points.}
The points $\sqrt{0}$, $\mathrm{arcosh}(1)$, $\mathrm{hypot}(0, 0)$ all
share the form \emph{the natural EML witness collides with the
junk-value boundary of $\log$}. The paper itself does not provide
explicit EML terms for these; we document them as out-of-scope
boundary artefacts, with concrete counterexample evaluations.

\paragraph{Domain-narrowing companions.}
The Euler bridges as stated narrow at the principal-branch cut. To
widen, we introduce \emph{negative-side companions}: a witness
$\mathsf{T}_f^{-}$ for the negative half-axis, constructed by feeding
$-x$ \textemdash\ the real EL expression $\mathrm{neg}(\mathbf{x}_0)$
lifted along $\iota$ \textemdash\ into the same Euler-bridge skeleton with
appropriately flipped constants. The two witnesses (positive and
negative) seal a substantial open neighbourhood of $0$, with $x = 0$
itself filled in by a trivial constant witness.

% =====================================================================
\section{A worked example: \texorpdfstring{$\sin x$}{sin x} end-to-end}
% =====================================================================

We trace $\sin$ through every architectural layer, as the canonical
case combining all the above.

\paragraph{Step 1 \textemdash\ source.} The paper claim is
\[
  \exists\,\mathsf{T}_{\sin} \in \EsetC,\;
  \forall\, x \in (0, \pi),\;
  \exists\, v \in \mathbb{C},\;
  \mathsf{T}_{\sin} \downarrow v
  \;\wedge\; \Re(v) = \sin(x).
\]

\paragraph{Step 2 \textemdash\ choose the witness.} From the Euler list
\eqref{eq:euler}, take the entry for $\sin$ and assemble its
$\EsetC$ realisation:
\[
  \boxed{\;\;\mathsf{T}_{\sin} \;:=\;
    \hexpC\!\bigl(\;
      \hlogC(\mathsf{C}_{\cos}) \;\hsubC\; \hlogC(\mathsf{I})
    \;\bigr).\;\;}
\]
Here $\mathsf{C}_{\cos}$ is the cosine-direction witness whose
evaluation is $e^{ix}$ for $x > 0$; it is itself a closed-constants +
$\mathbf{x}_0$ assembly using $\mathsf{I}$ and $\haddC$ to encode
$\log i + \log x = \log(ix)$ before two outer $\hexpC$'s.

\paragraph{Step 3 \textemdash\ evaluation chain.} Five sub-steps, each
one closure-lemma application:
\begin{enumerate}[label=(\roman*), topsep=2pt, itemsep=2pt]
  \item $\mathsf{C}_{\cos} \downarrow e^{ix}$ \quad (Euler bridge for
    $e^{ix}$, itself a small instance of the same template).
  \item $\hlogC(\mathsf{C}_{\cos}) \downarrow ix$ \quad (closure for
    $\hlogC$; side condition $\Im(ix) = x \in (-\pi, \pi]$ \checkmark).
  \item $\hlogC(\mathsf{I}) \downarrow i\pi/2$ \quad (closure for
    $\hlogC$, plus the Mathlib value $\log i = i\pi/2$).
  \item $\hlogC(\mathsf{C}_{\cos}) \hsubC \hlogC(\mathsf{I})
    \downarrow i(x - \pi/2)$ \quad (closure for $\hsubC$).
  \item $\mathsf{T}_{\sin} \downarrow \exp\!\bigl(i(x - \pi/2)\bigr)$
    \quad (closure for $\hexpC$, unconditional).
\end{enumerate}

\paragraph{Step 4 \textemdash\ projection.} Take the real part:
\[
  \Re\!\Bigl(\exp(i(x - \tfrac{\pi}{2}))\Bigr)
   = \cos\!\Bigl(x - \tfrac{\pi}{2}\Bigr)
   = \sin x,
\]
using the cofunction identity $\cos(x - \pi/2) = \sin x$. \(\square\)

\medskip

Every layer of the architecture appears here in concentrated form:
the grammar of \S3, the catalogue of \S\ref{sec:combinators}, the
closed constant $\mathsf{I}$ and the Euler closed-constant scaffolding
of \S\ref{sec:euler}, the real-to-complex lift implicit in any use of
$\hsubC$, the Euler identity for $\sin$, and the principal-branch
projection. The fact that the same template handles all six trig
primitives \textemdash\ with only the inner content of $\mathsf{C}_f$ and
the choice of projection $\Re$ or $\Im$ varying from line to line of
\eqref{eq:euler} \textemdash\ is what gives the proof its uniformity.

% =====================================================================
\section{What is novel}
% =====================================================================

\begin{enumerate}[topsep=2pt, itemsep=4pt]
\item \textbf{The structural compiler.} Theorem~\ref{thm:compile} is a
  single induction with one combinator application per constructor.
  Each elementary composition reduces \emph{uniformly} \textemdash\ no
  per-primitive cleverness. The heavy lifting (the eleven-step
  unfolding for $\hadd$, the deep unfolding for $\hsqrt$, etc.) is
  amortised inside the catalogue, proved once and re-used.

\item \textbf{Compositional branch-cut bookkeeping.} The closure
  lemmas for $\haddC$ and friends carry an explicit eleven-field
  precondition (the so-called \emph{ADDsafe} bundle) that summarises
  all imaginary-part-in-strip and non-vanishing conditions the
  unfolding generates. Every trig witness's evaluation proof reduces
  to discharging this precondition by a few line-of-reasoning
  numerical inequalities. The branch-cut handling becomes
  \emph{compositional}, not ad-hoc per primitive.

\item \textbf{The real-to-complex lift.} The homomorphism
  $\iota : \EML \to \EMLC$ lets complex trig witnesses re-use real
  arithmetic that has already been structurally compiled. Without
  $\iota$, every trig witness would have to reprove $\sqrt{1 - x^2}$
  from first principles in the complex grammar; with it, the witness
  for $\arccos$ literally uses the real $\sqrt{1 - x^2}$ tree, lifted.

\item \textbf{Negative-side companions, from one observation.}
  A surprising amount of trig domain expansion ($\sin$ to $(-\pi, \pi)$,
  $\cos$ to $\mathbb{R}\setminus\{0\}$, $\arctan$ to
  $(-\pi, \pi)\setminus\{0\}$, etc.) follows from a single observation:
  the universal blocker is the same side condition (the strict
  inequality $\arg < \pi$ for $\hlogC$). Once $-x$ is encoded via
  $\iota$, the same Euler-bridge skeleton runs on the flipped argument
  and produces the cofunction-shifted target. One template, five
  widenings.
\end{enumerate}

% =====================================================================
\section{The role of the principal branch}
\label{sec:branch}
% =====================================================================

The fixed choice of branch for the complex logarithm is not a cosmetic
detail of the construction; it is structural. Almost every quantitative
constraint in the proof \textemdash\ the narrowed trig domains, the
side-condition bookkeeping in the closure lemmas, the negative-side
companion witnesses, and the structural exclusion of certain boundary
points \textemdash\ traces back to the principal-branch convention.
This section unpacks how.

\paragraph{The convention.}
Throughout the construction we use the principal branch
\[
  \log z \;=\; \log|z| \,+\, i\,\arg(z),
  \qquad \arg(z) \in (-\pi, \pi].
\]
This is the convention of essentially every standard mathematical
library. Its critical feature for our purposes is the inversion
identity
\begin{equation}
  \log(\exp w) \;=\; w
  \qquad\text{provided}\qquad
  \Im(w) \;\in\; (-\pi, \pi].
  \label{eq:strip}
\end{equation}
Outside this strip, $\log \circ \exp$ wraps by integer multiples of
$2\pi i$, breaking the identity.

\paragraph{How the constraint enters the closure lemmas.}
Recall the logarithm combinator
\[
  \widehat{\log}(T) \;:=\;
    \mathrm{eml}\!\bigl(\mathbf{1},\,
      \mathrm{eml}(\mathrm{eml}(\mathbf{1}, T),\, \mathbf{1})\bigr).
\]
For $T \downarrow v$, the chain of $\mathrm{eml}$-unfoldings produces
intermediate values whose imaginary parts must lie in the strip
\eqref{eq:strip} for the inversion to fire. The decisive step is the
outermost $\mathrm{eml}(\mathbf{1}, \cdots)$, where one applies
$\log\!\circ\!\exp$ to the value $e - \log v$, which has imaginary part
\[
  \Im(e - \log v) \;=\; -\,\Im(\log v) \;=\; -\,\arg(v).
\]
The strip condition $\Im(\cdot) \in (-\pi, \pi]$ thus reads
$-\arg(v) \in (-\pi, \pi]$, equivalently
\[
  \arg(v) \;\in\; [-\pi, \pi),
\]
which excludes only one possible value: $\arg(v) = \pi$, the negative
real axis. The closure lemma for $\widehat{\log}$ therefore carries
the side condition
\[
  v \ne 0 \quad\text{and}\quad \arg(v) < \pi \;\text{strictly}.
\]
\emph{This is the core branch-cut constraint of the entire proof.}

\paragraph{The cascade.}
Every higher combinator that internally applies $\widehat{\log}$
inherits this constraint. In particular:
\begin{itemize}[topsep=2pt, itemsep=2pt]
  \item $\hmulC$ is implemented as
    $\hexp\!\bigl(\hlogC(A) \haddC \hlogC(B)\bigr)$, so it requires
    \emph{both} $A.\!\!\downarrow$ and $B.\!\!\downarrow$ to satisfy the
    branch condition.
  \item $\hdivC$ is implemented analogously and inherits the same
    constraints from each operand.
  \item Witnesses constructed by composition (the trig family, the
    inverse trig family, $\sqrt{\,}$ in the complex extension) carry
    a sum of constraints \textemdash\ one per occurrence of $\hlogC$
    inside the witness.
\end{itemize}
This cascade is what determines the natural sealed subdomain of every
trig primitive.

\paragraph{Concrete consequences.}
Three of the most visible narrowings have a direct branch-cut origin.
\begin{itemize}[topsep=2pt, itemsep=4pt]
  \item \emph{$\sin$ on $(0, \pi)$.} The witness applies $\hlogC$ to
    $e^{ix}$, whose argument is $x$; the side condition $\arg < \pi$
    becomes $x < \pi$. Pushing past $\pi$ would put $\arg(e^{ix})$ on
    the cut.
  \item \emph{$\arcsin$ originally on $(0, 1)$.} The witness multiplies
    $i$ by $x$ via $\hmulC$, which requires $\arg(x) < \pi$;
    \emph{positive} reals have $\arg = 0$ and pass, but \emph{negative}
    reals have $\arg = \pi$ and fail. The bound $0 < x$ on the original
    witness is not a topological accident \textemdash\ it is exactly
    "$x$ is not on the branch cut".
  \item \emph{$\tan$ on $(0, \pi/2)$.} The Cayley quotient witness
    invokes $\hmulC$ on $2x$, again hitting the same constraint.
\end{itemize}

\paragraph{Branch-cut workarounds, made systematic.}
The widening companions of \S7 are best understood as a unified
collection of \emph{branch-cut workarounds}, all leveraging the same
observation: \emph{if the natural witness fails because some
sub-expression is a negative real, replace that sub-expression with
$-(-\,\bullet)$, lifted into the complex grammar via $\iota$, so the
inner $\bullet$ is positive}.

For each of the five widenings, this reduces to choosing a different
analytic identity:
\begin{itemize}[topsep=2pt, itemsep=2pt]
  \item For $\arcsin$, use $\arcsin x = \tfrac{\pi}{2} - \arccos x$.
    The witness for $\arccos$ already covers $(-1, 1)$ in full because
    its inner multiplication is $i \cdot \sqrt{1 - x^2}$, with
    $\sqrt{1-x^2} > 0$ on the open interval, dodging the cut entirely.
  \item For $\arctan$ on the negative side, use $1 + ix = 1 - i\cdot(-x)$
    with $-x > 0$.
  \item For $\cos$, use the parity $\cos(-x) = \cos(x)$ to feed $-x$
    into the same Euler skeleton.
  \item For $\sin$, use $\sin x = \cos(\tfrac{\pi}{2} - x)$ together
    with $\log(-i) = -i\pi/2$ to keep the projection on the real part.
  \item For $\tan$ on the negative side, use the swap-numerator Cayley
    quotient $(1 - e^{-2ix}) / (1 + e^{-2ix}) = i\tan x$.
\end{itemize}
Each identity is the analytic content; the branch cut is the technical
constraint each identity is chosen to dodge.

\paragraph{An alternative would change the entire architecture.}
The paper itself sketches a different approach: redefine the branch
of $\log$ inside the EML evaluator so that the cut runs elsewhere.
Concretely, the paper's compiler "manually corrects" the sign of $i$
when crossing the cut, effectively choosing a branch matched to the
specific witness shapes used. With such a tailored branch, several
trig primitives could be sealed on their full real domain by a single
witness rather than by positive- and negative-side companions.

The price is structural: each closure lemma in the catalogue would
have to be reproved against the alternative branch. Our construction
chooses the standard convention \eqref{eq:strip} and pays the price
in narrowed domains. Both choices are mathematically valid; they
trade simplicity of side conditions for breadth of the sealed
subdomain.

% =====================================================================
\section{Open questions}
\label{sec:open}
% =====================================================================

The construction above resolves the headline existence claim of
Theorem~\ref{thm:headline}. Several structural and quantitative
questions remain. We organise them into three groups: questions
inherited from the paper's own list, questions about strengthening
the present construction, and questions about the optimal expression
sizes.

\subsection*{Questions from the paper}

The paper's Supplementary Information lists seven open problems that
neither the paper nor the present construction resolves. We restate
the most accessible ones.

\begin{enumerate}[topsep=2pt, itemsep=4pt, label=\textbf{P\arabic*}]
\item \textbf{Constant-free binary Sheffer.}
  Does there exist a binary operator $\star : \mathbb{R}^2 \to \mathbb{R}$
  (or $\mathbb{C}^2 \to \mathbb{C}$) such that $\{\mathbf{x}_n, \star\}$
  alone, with no distinguished constant, generates all elementary
  functions? An exhaustive numerical search to operator complexity
  $K = 6$ has found no candidate, but no proof of impossibility is
  known.

\item \textbf{Real-only Sheffer.}
  Does there exist a binary operator on $\mathbb{R}$ (no complex
  intermediates) that forms a complete basis for the elementary
  functions? The paper conjectures the answer is \emph{no}: every
  candidate appears to require either a branch through $\mathbb{C}$
  to compute $\pi$ and the trigonometric family, or an extended-real
  intermediary $-\infty$. A formal impossibility proof is open.

\item \textbf{$-\infty$-free EML.}
  Can the EML operator (or one of its cousins EDL, $-\textsc{eml}$) be
  used without the extended real axis, in particular without
  $-\infty$ as a possible intermediate value? In our partial-evaluation
  presentation we side-step this by returning $\bot$ instead of
  $-\infty$, but the question of whether the elementary functions are
  expressible \emph{purely on the real line} is independent.

\item \textbf{Universal minimality of $\{1, \mathrm{eml}\}$.}
  The pair (one constant, one binary operator) is empirically the
  smallest known calculator for the elementary functions. A truly
  universal minimality theorem would quantify over every continuous
  binary operator and every constant, asserting that no two-element
  basis is strictly smaller in any reasonable measure. The paper gives
  a cautionary trap example, $B(x, y) = x - y/2$, to illustrate why
  naive arguments fail.

\item \textbf{Stern--Brocot enumeration for EML expressions.}
  The set of EML expressions is enumerable, but with massive
  redundancy: many syntactically distinct trees evaluate to the same
  elementary function. Is there a canonical form, analogous to the
  Stern--Brocot tree for rationals, that enumerates each elementary
  function exactly once?

\item \textbf{Taxonomy of the Sheffer family.}
  The known cousins are EML, EDL ($\exp(x)/\log(y)$, requiring the
  constant $e$), and $-\textsc{eml}$ ($\log(x) - \exp(y)$, requiring
  $-\infty$). Are these three points isolated, members of a discrete
  family, or random samples of a continuous one-parameter family of
  Sheffer operators?
\end{enumerate}

\subsection*{Questions about the present construction}

\begin{enumerate}[topsep=2pt, itemsep=4pt, label=\textbf{C\arabic*}]
\item \textbf{Full-real-domain trig.}
  Our construction seals the trig family on wide symmetric subdomains
  around $0$, using positive- and negative-side companions. Can the
  same family be sealed on the \emph{full} real domain by a single
  witness? Two paths suggest themselves:
  \begin{itemize}[topsep=2pt, itemsep=2pt]
    \item \emph{Via a custom branch}: re-derive the catalogue against
      the paper's "manually corrected" branch. The closure lemmas
      change, but the witnesses gain wider sealed subdomains by
      construction.
    \item \emph{Via multi-witness periodicity}: a witness family
      indexed by period number, with each member covering one
      fundamental period. The existential becomes weaker
      ($\forall x.\,\exists\,T$ rather than $\exists\,T.\,\forall x$),
      but the construction stays inside the principal branch.
  \end{itemize}
  Both are tractable but require non-trivial new infrastructure.

\item \textbf{Per-primitive completeness for the cousins.}
  Theorem~\ref{thm:headline} concerns the EML operator. The paper
  presents EDL and $-\textsc{eml}$ as confirmed-empirically-equivalent
  cousins; both deserve a parallel completeness theorem. The grammar
  scaffolding is straightforward (replace $\mathrm{eml}$ by the cousin
  operator); the witness search is genuinely new work, since the
  cousin operators have a different algebraic shape and the EML
  witnesses do not directly translate.

\item \textbf{The three structural boundary points.}
  $\sqrt{0}$, $\mathrm{arcosh}(1)$, $\mathrm{hypot}(0, 0)$. Each fails
  for the same reason: the natural EML expression for the value
  collides with the convention $\log(0) = 0$ that any total
  formalisation must adopt. Two candidate fixes:
  \begin{itemize}[topsep=2pt, itemsep=2pt]
    \item Extend the grammar with a primitive $\mathrm{rpow}$ operator
      that bypasses the $\log$ at the boundary.
    \item Move the affected witnesses entirely into the complex
      extension, where the boundary lives in different coordinates.
  \end{itemize}
  Neither has been pursued; both would extend the grammar beyond
  the paper's pure $\{1, \mathrm{eml}\}$ alphabet.
\end{enumerate}

\subsection*{Quantitative questions on tree size}

\begin{enumerate}[topsep=2pt, itemsep=4pt, label=\textbf{S\arabic*}]
\item \textbf{Optimal tree size per primitive.}
  Our construction produces witnesses whose tree sizes range from
  $K = 7$ (for $0$) to $K \approx 10^{7}$ (for $\log_b$). The paper
  records small hand-tuned figures (e.g.\ $K = 5$ for $\sin$ via the
  identity $\sin x = \cos(x - \pi/2)$ in a setting where $\cos$ is
  already available). The gap reflects that our witnesses are
  produced by a uniform structural compiler, not optimised per
  primitive. What is the smallest $K$ achievable for each primitive
  under the principal-branch convention?

\item \textbf{Asymptotic depth.}
  The paper observes that the identity function appears at depth $4$
  in the EML tree, allowing variable transplanting by multiples of
  $4$. Are there non-trivial functions whose minimum-depth witness has
  a different residue class modulo $4$? If so, what depths are
  achievable?

\item \textbf{Compositional bounds.}
  Given $K(f), K(g)$, what is the best upper bound on $K(f \circ g)$
  derivable from the EML grammar? Our compiler achieves something like
  a multiplicative bound; tighter additive or near-additive bounds
  would substantially shrink the witnesses for composite primitives
  (and thereby narrow the gap with the paper's hand-tuned figures
  in S1).
\end{enumerate}

% =====================================================================
\section{Coverage scoreboard}
% =====================================================================

\begin{center}
\begin{tabular}{l l l}
\toprule
\textbf{Family} & \textbf{Sealed coverage (post-submission)} & \textbf{Mechanism} \\
\midrule
Atoms (7)              & full                                                & direct closed witness \\
Real unary (8)         & full $\setminus \{\sqrt{0}\}$                       & structural compiler \\
Real binary (8)        & full $\setminus \{\mathrm{hypot}(0,0)\}$            & structural compiler \\
Hyperbolic (6)         & full $\setminus \{\mathrm{arcosh}(1)\}$             & $\F\to\EL$ + compiler \\
$\cos$, $\arccos$, $\arcsin$ (3) & full natural domain                       & Euler bridges + companions \\
$\sin$, $\arctan$, $\tan$ (3) & full natural domain (Path~C$'$)             & Euler bridges + range-reduction \\
\bottomrule
\end{tabular}
\end{center}

\medskip

\noindent\textbf{Net.} $36/36$ paper primitives sealed by finite EML
expressions on their natural domains (modulo three structurally
excluded $\sqrt{0}$-collision boundary points). The trig family was
narrowed to symmetric subdomains around $0$ at submission time; the
post-submission Path~C$'$ widening (\S\ref{sec:post-submission}) lifts
that narrowing for $\sin$, $\arctan$, $\tan$. The headline theorem
(Theorem~\ref{thm:headline}) holds in its strongest form.

\medskip

\noindent\textbf{Public Lean API.} $48$ \texttt{paper\_claim\_*}
theorems exported by \texttt{EML.Framework.PaperClaims}, plus $8$
\texttt{edl\_paper\_claim\_*} (Plan~D) and $5$
\texttt{negEml\_paper\_claim\_*} (Plan~E, including the
\texttt{minusInf} witness via the parallel
\texttt{NegEMLTermE} grammar over \texttt{EReal}) in
\texttt{EML.Framework.Sheffer} \textemdash\ a total of $61$
machine-checked completeness theorems. Local \texttt{lake build}:
$8056$ jobs, sorry-free. Independent re-verification on PCSS Eagle
HPC: clean.

% =====================================================================
\section{Addendum: post-submission progress (Path~C$'$ + Plans A/D/E)}
\label{sec:post-submission}
% =====================================================================

This section records four developments after the submission deck froze.
The mathematical content of the original construction is unchanged;
these are extensions, cleanups, and Sheffer-cousin progress.

\subsection*{Plan~A \textemdash\ Sheffer naming cleanup}

The submission-era \texttt{Sheffer.lean} contained four operator
grammars (\texttt{EDLTerm}, \texttt{LDETerm},
\texttt{T1Term}, \texttt{T2Term}). Three of those names were misnomers
relative to the paper's \S3.1:
\texttt{LDE}~$=\log(x)/\exp(y)$ (division) was misnamed as the paper's
\textbf{$-$EML}~$=\log(x)-\exp(y)$ (subtraction); \texttt{T1Term} and
\texttt{T2Term} were binary, but the paper's actual T$_1$ and T$_2$ are
\emph{ternary} operators (Supplementary Information \S1.4) flagged by
the paper itself as preliminary unverified candidates for a constant-free
Sheffer.

After cleanup, \texttt{Sheffer.lean} hosts exactly the two
paper-named cousins, $\mathrm{EDL}$ and $-\mathrm{EML}$, with the
correct definitions and a line-level paper-sourcing audit trail in
\texttt{notes/legacy\_planning/Sheffer\_PaperSourcing.md}.
The fabricated binary T$_1$/T$_2$
are removed; the paper's real ternary T$_1$/T$_2$ are documented as
preliminary future work, not as part of the present scaffolding.

\subsection*{Plan~B \textemdash\ architectural finding (custom log
branch infeasible)}

The paper at line~333 sketches an alternative route to full real-domain
trig coverage by ``manually correcting the $i$-sign'' inside the
compiler. Implementing this in Lean was the original Plan~B and
foundered on a structural fact: the $\EML_{\mathbb{C}}$ partial
evaluation (\texttt{EMLTerm}\textsuperscript{$\mathbb{C}$}.\texttt{eval?}
in \texttt{Framework/Complex/Term.lean}) hard-codes Mathlib's
principal-branch \texttt{Complex.log}; there is no place in the
grammar to substitute a different branch convention.

The constructive observation \textemdash\ that \texttt{mkLog}$_{\mathbb{C}}T$
\emph{does} evaluate at the cut $\arg v = \pi$, returning
$\mathrm{Complex.log}\,v - 2\pi i$ instead of
$\mathrm{Complex.log}\,v$, so any witness whose final operation is
$\mathrm{exp}_{\mathbb{C}}$ (such as $\cos$) absorbs the
$2\pi i$-shift via $\exp$-periodicity \textemdash\ identifies why $\cos$
already covers $\mathbb{R}\setminus\{0\}$ but $\sin$, $\arctan$, $\tan$
do not. The fix-point of this analysis is that the paper's
``manual correction'' is necessarily a meta-level operation choosing a
different witness shape per region of $x$, mathematically equivalent
to the witness-family construction of Path~C$'$.

\subsection*{Plan~C$'$ \textemdash\ full-real-domain trig via
range-reduction by substitution}

A GPT~Pro consultation (transcript in \texttt{gpt\_pro\_bundle/}) gave
a refined plan that uses the existing positive/negative companion
witnesses but lifts them to full real domain by \emph{substitution}.
Three engineering moves:

\begin{enumerate}[topsep=2pt, itemsep=2pt]
  \item \textbf{Variable substitution on $\EsetC$.} Define
    $\mathrm{subst}_0 : \EsetC \to \EsetC \to \EsetC$ that replaces
    every $\mathbf{x}_0$ in $t$ with a given term $s$. Prove
    $\mathrm{eval}_{?}(t.\mathrm{subst}_0\,s,\,\mathrm{env})
      = \mathrm{eval}_{?}(t,\,\mathrm{env}[\mathbf{x}_0 \mapsto
      \mathrm{eval}_{?}(s,\,\mathrm{env})])$. (\texttt{Subst.lean}.)

  \item \textbf{Real-safe addition.} The lemma
    $\mathrm{ADDsafe}_{\mathbb{C}}((\overline{a},0), (\overline{b},0))$
    \textemdash\ the gnarly 11-condition $\mathrm{mk}\widehat{+}_{\mathbb{C}}$
    precondition bundle \textemdash\ holds automatically when both
    arguments are real-valued, the only nontrivial case being
    $\exp(a) - a \neq 0$ (true since $\exp(a) \geq a + 1$). With this
    foundation lemma, every period-shift via repeated
    $\mathrm{mk}\widehat{+}_{\mathbb{C}}$ stays in the real fragment,
    so the $\arg = \pi$ boundary trap never appears.

  \item \textbf{Reuse, don't reshape.} For each narrow primitive,
    pick an algebraic identity that reduces to a primitive whose
    witness is \emph{already} full-domain or full-natural-open:
    \begin{itemize}[topsep=2pt, itemsep=2pt]
      \item $\sin x = \cos(\pi/2 - x)$: substitute $\pi/2 - x$ for
        $\mathbf{x}_0$ in the existing $\cos$ witness; finish with
        \texttt{Real.cos\_pi\_div\_two\_sub}. Coverage:
        $\mathbb{R}\setminus\{\pi/2\}$.
      \item $\arctan x = \arcsin(x/\sqrt{1+x^2})$: substitute the
        real-fragment compile of $x/\sqrt{1+x^2}$ for $\mathbf{x}_0$
        in $\arcsin$'s already-full witness; finish with
        \texttt{Real.arctan\_eq\_arcsin}. Coverage: full $\mathbb{R}$.
      \item $\tan x = \tan(x - k\pi)$: range-reduce arbitrary $x$ with
        $\cos x \neq 0$ to the fundamental strip $(-\pi/2, \pi/2)$ via
        repeated $\pm\pi$ shifts; apply the existing local Cayley
        witness on the reduced argument; finish with
        \texttt{Real.tan\_sub\_int\_mul\_pi}. Coverage:
        $\{x : \cos x \neq 0\}$.
    \end{itemize}
\end{enumerate}

\noindent\textbf{Result.} Three new public theorems
$\mathrm{paper\_claim}_{\sin\_\mathrm{full}}$,
$\mathrm{paper\_claim}_{\arctan\_\mathrm{full}}$,
$\mathrm{paper\_claim}_{\tan\_\mathrm{full}}$ \textemdash\
\emph{witness-family} statements of the shape
\[
  \forall\, x \in \mathrm{Domain}(f),\;
  \exists\, t \in \EsetC,\;
  \exists\, v \in \mathbb{C},\;
  t \downarrow v \;\wedge\; \pi(v) = f(x).
\]
The witness depends on $x$'s region (positive/negative side, or
period number for $\tan$); paper-faithfulness is preserved via the
analogy with the paper compiler's manual sign correction.

\subsection*{Plan~D \textemdash\ EDL per-primitive completeness}

The paper's \S3.1 promotes the EDL cousin
$\mathrm{edl}(x, y) = \exp(x) / \log(y)$ to the same completeness
status as EML, but supplies no Lean-style proof. Eight of the $36$
paper primitives are now sealed in our framework via Aristotle-aided
witness search:

\begin{center}
\begin{tabular}{l l l}
\toprule
\textbf{Primitive} & \textbf{Witness} & \textbf{Provenance} \\
\midrule
$1$, $\mathbf{x}_0$, $e$       & atomic                                                         & grammar \\
$\exp x$                       & $\mathrm{edl}(\mathbf{x}_0, e)$                                & identity (collapse) \\
$\log x$                       & $\mathrm{edl}(1, \mathrm{edl}(\mathrm{edl}(1, \mathbf{x}_0), e))$ & Aristotle \\
$x / y$                        & $\mathrm{edl}(\mathrm{D8}(\mathbf{x}_0), \mathrm{D4}(\mathbf{x}_1))$ & Aristotle \\
$\exp(\exp x)$                 & $\mathrm{edl}(\mathrm{edl}(\mathbf{x}_0, e), e)$              & Aristotle \\
$\log(\log x)$                 & nested D8                                                       & Aristotle \\
\bottomrule
\end{tabular}
\end{center}

\noindent\textbf{Structural ceiling.} Aristotle's analytical work
(transcripts in \texttt{chunks/085\_*}, \texttt{chunks/089\_*}) shows
that the EDL combinator
$\mathrm{edl}(a, b) = \exp(a)/\log(b)$ admits no mechanism for
\emph{addition} of sub-expression values. Multiplication ($x \cdot y =
\exp(\log x + \log y)$), squaring, square-root, and the entire
trigonometric and hyperbolic family are therefore conjecturally
unreachable from closed EDL terms over $\{1, e, \mathrm{var},
\mathrm{edl}\}$. The constants $-1$, $2$, $\frac12$ are also conjecturally
unreachable on Schanuel-conjecture grounds (their construction would
require $\log 2 \in$ EL-closure of $\{1, e\}$, which Schanuel forbids).
This validates the paper's own framing of EDL completeness as
\emph{conjectured}, not proven.

\subsection*{Plan~E \textemdash\ $-$EML pilot}

The third Sheffer cousin in the paper's \S3.1 is $-$EML, paired with
the constant $-\infty$. Two trivial atoms ($1$, $\mathrm{var}$) are
sealed in our $\mathbb{R}$-based \texttt{NegEMLTerm}. Sealing the
$-\infty$ constant requires switching the grammar to $\mathrm{EReal}$
(extended reals); a self-contained pilot in
\texttt{chunks/088\_neg\_eml\_pilot} demonstrates the approach.
Per-primitive completeness for $-$EML is paper-open, with the same
addition-unreachability obstruction as Plan~D applying to the
arithmetic and trig families.

\subsection*{Aristotle integration}

Of the 10 chunks submitted in this round, 9 returned successful
proofs, accounting for: two pure-Mathlib auxiliaries
(\texttt{atanArg\_in\_Ioo}, \texttt{tan\_period\_reduction}); three
end-of-pipeline assembly proofs for Path~C$'$
($\mathrm{sin\_via\_cos}$, $\mathrm{arctan\_via\_arcsin}$,
$\mathrm{tan\_full}$); five Plan~D / Plan~E witness searches
(D1\textendash D11, E1\textendash E3). Aristotle's role was strictly
within the \emph{audit blade} category: the architectural design
came from human + GPT~Pro, the framework integration was
hand-coded, and Aristotle's contribution was witness search and
proof-tactic discovery within self-contained chunk files.

\end{document}
